\clearpage\thispagestyle{otFront}\phantomsection\label{avant-propos}\pdfbookmark[0]{Apprendre et comprendre}{bm-avant-propos}\otPageTitre{Apprendre et comprendre}{Un livre pour toute l’année de Troisième}\fontsize{11}{15}\selectfont

Ce livre accompagne la classe de Troisième et la préparation du brevet 2029,
avec le nouveau programme de mathématiques du cycle 4. Il aide à consolider les
calculs, à comprendre les nouvelles notions, puis à résoudre des problèmes en
reliant plusieurs outils. Il complète le travail de classe, dans l’ordre adapté aux besoins.\par\medskip
\begin{savoirfaire}[Une progression en trois temps]Chaque chapitre propose deux pages de cours, cinq exercices de découverte, cinq d’entraînement et cinq problèmes. Les corrections explicitent la méthode, les calculs et l’interprétation. Une figure utile est reprise avec le corrigé ; une construction cherchée n’apparaît qu’à cet endroit.\end{savoirfaire}
\begin{savoirfaire}[Chercher avant de consulter]Écrire une première démarche même incomplète. Après correction, fermer le livre et refaire les étapes décisives. Quelques jours plus tard, reprendre un exercice proche sans aide : reconnaître un résultat n’est pas encore savoir le retrouver.\end{savoirfaire}
\begin{savoirfaire}[Aller plus loin]Les dix-huit défis demandent une initiative supplémentaire, avec deux indications séparées et une solution rédigée. Deux brevets blancs permettent un entraînement chronométré. Les fiches et la matrice finale aident à repérer la notion à revoir.\end{savoirfaire}
Les six compétences travaillées sont chercher, modéliser, représenter,
raisonner, calculer et communiquer. Un dessin, un tableau ou un script sont
des moyens de penser autant que des supports de réponse.

\clearpage
